% ======================================================================
% Drop-in LaTeX drafts answering the four review risks (2026-07-24).
% Every number below is measured from the repo (sources in comments).
% ======================================================================

% ----------------------------------------------------------------------
% 1) REPRODUCIBILITY ADDENDUM  (append to Sec. Simulation Settings)
% Facts: sim/real_fl.py (_prep_data, RealMFL.__init__, local_train,
% _batch), sim/face.py (_adopt), sim/multimodal_model.py (FusionHead),
% sim/run_v2x_real.py (round loop), sim/config.py.
% [FIX] The current settings text says vehicles hold "80--600 local
%       training objects": the code gives ~83 (data-rich) and 4--11
%       (all others) on KITTI. Replace that sentence with the below.
% [VERIFY] the nuScenes rich-vehicle share before quoting its number.
% ----------------------------------------------------------------------
\textbf{Data partition and training protocol.}
Both datasets are class-balanced by downsampling every class to the
smallest kept class (KITTI: $1{,}437$ objects per class over Car /
Pedestrian / Cyclist; nuScenes: Car / Pedestrian after dropping classes
with fewer than $800$ samples) and split $68/12/20\%$ into
train/validation/test by a seeded permutation; validation and test sets
are shared across vehicles and never corrupted. Each vehicle draws a
disjoint local set from the training pool: the $15\%$ data-rich
vehicles split the residual pool equally (${\approx}83$ objects each on
KITTI), while every remaining vehicle holds only $4$--$11$ objects,
freezes its encoders, and trains only its lightweight fusion head --
its accuracy is therefore determined by the encoders it receives.
Label distributions are uniform by construction; heterogeneity enters
through data volume and sensing quality: data-rich vehicles draw a
quality factor $Q\sim\mathcal U(0.8,1)$, the rest
$Q\sim\mathcal U(0.1,0.3)$, and low-quality vehicles train on corrupted
inputs (camera: additive noise $\sigma{=}0.25$ with $0.6\times$
brightness; LiDAR: $40\%$ point dropout). Encoders are
modality-specific networks producing fixed-width features, fused by a
two-layer MLP head (concatenation $\to 64 \to$ classes). Each round
executes (i)~one local epoch of Adam (lr $10^{-3}$, batch $\le 64$);
(ii)~publication and V2V exchange (Sec.~III); (iii)~evaluation-gated
adoption: each vehicle ranks received candidates by the optimistic gain
estimate, evaluates the top $N_{\mathrm{ev}}{=}6$ on the shared
validation split by temporarily swapping in the candidate encoder, and
adopts the best iff the measured validation gain reaches
$\delta{=}5\times10^{-3}$, merging via data-size-weighted
FedAvg~\cite{fedavg}; (iv)~test evaluation. Results average three runs
(seeds varying data assignment, modality availability, and matching);
$\pm$ denotes the standard deviation over runs.

% ----------------------------------------------------------------------
% 2) HINDSIGHT-ORACLE COMPARISON  (add to Sec. Evaluation; discharges
%    the twice-promised clairvoyant comparison)
% Data: results/face_oracle.npz (sim/face_oracle.py, 16 ILP instances);
% figure panel ready in sim/face_figs.py (~line 221).
% ----------------------------------------------------------------------
\textbf{Comparison with a clairvoyant oracle.}
To position FACE against the best possible forwarding decisions, we
extract $16$ small sub-instances ($18$-vehicle neighborhoods over short
windows) of the Seoul trace and solve each offline as an integer
program given full knowledge of all future contacts, link budgets, and
realized encoder gains. Running causally with no future information,
FACE attains $0.53\times$ the clairvoyant objective on average (median
$0.60\times$), i.e., within a factor of two of hindsight-optimal. A
clairvoyant variant with relaying disabled attains $0.96\times$,
showing that on instances this small the oracle's advantage stems
almost entirely from privileged scheduling knowledge rather than from
any mechanism FACE lacks.
% (caveat for rebuttals: relay value grows with instance size; the
% 18-vehicle windows fit the ILP, so this understates ferrying gains.)

% ----------------------------------------------------------------------
% 3) SURROGATE FIDELITY  (reframe the theorem; option 2 of the review)
% Data: newnewdata/theory_stats.txt --
%   eps_pred over n=20,943 adoption evaluations (Table-I runs):
%     mean |pred-realized| = 0.0077, p95 = 0.0196,
%     mean realized gain 0.0039, mean predicted 0.0043 (bias +11%).
%   eps_int over n=500 concurrent-delivery probes: 0.0575, p95 0.386.
% Also: the manuscript references subsec:agg_train, which does not
% exist -- restore the subsection or fix the label.
% ----------------------------------------------------------------------
We state the surrogate result as a \emph{fidelity lemma}, not a
performance guarantee: it bounds the surrogate-to-realized-utility gap
in terms of the prediction and interaction errors
$\bar\epsilon_{\mathrm{pred}},\bar\epsilon_{\mathrm{int}}$, which we
measure rather than assume. Over the $20{,}943$ adoption evaluations of
the main runs, the per-event absolute prediction error is
$\bar\epsilon_{\mathrm{pred}}{=}0.008$ (p95 $0.020$) against a mean
realized gain of $0.004$; a worst-case accumulation over
$T|\mathcal I|$ terms is therefore vacuous, and we do not claim it.
What the deployment relies on -- and what the data support -- is that
the errors are unbiased and rank-consistent: the signed bias is only
$+11\%$ of the mean per-event gain, and the decile-level calibration of
predicted against realized gains is monotone
(Fig.~\ref{fig:calib}), so prediction errors cancel across events
instead of compounding, and the measured interaction term
($\bar\epsilon_{\mathrm{int}}{=}0.058$ over $500$ concurrent-delivery
probes) remains an order of magnitude below the per-round utility.

% ----------------------------------------------------------------------
% 4) BASELINE FAIRNESS + NEW BENCHMARKS  (extend Sec. Benchmarks)
% (a) budget-ordering clarification, (b) adaptation-scope honesty,
% (c) descriptions of the six new baselines now in sim/face.py
%     (results land in newnewdata/metrics_v2x_real_kitti_newbase.npz).
% ----------------------------------------------------------------------
% (a) append to the benchmark preamble:
When a contact budget cannot carry all offered encoders, every scheme
ranks its candidates by its own criterion and admits them greedily
under the same per-contact airtime knapsack and half-duplex matching as
FACE; demand-blind schemes (Cached-DFL, V2V) admit an arbitrary
feasible subset, matching their demand-blind semantics.
% (b) append after the mmFedMC/AutoFed items:
mmFedMC and AutoFed are server-based multimodal FL frameworks; we adapt
their \emph{selection principles} (contribution-per-communication-cost
and source-quality ranking, respectively) to the decentralized setting
and grant both the full Sec.~III protocol, so they differ from FACE
only in the forwarding value. Their server-side components (global
client sampling, cross-client imputation) have no decentralized
counterpart and are not ported.
% (c) additional benchmarks:
(6)~\textbf{PRoPHET}~\cite{prophet}: DTN routing adapted to encoder
forwarding -- per-pair delivery predictabilities with encounter
reinforcement, aging, and transitivity
($P_{\mathrm{init}}{=}0.75$, $\gamma{=}0.98$, $\beta{=}0.25$); an
encoder is forwarded to neighbors by their predictability of reaching a
current demander;
(7)~\textbf{Spray-and-Wait}~\cite{spraywait}: binary spray of the $K_x$
replication tickets (half the remaining tickets per encounter),
followed by a direct-delivery wait phase;
(8)~\textbf{Cached-DFL-LFU/-rand}: Cached-DFL with LFU / random cache
placement instead of LRU;
(9)~\textbf{No-comm} and \textbf{Full-contact}: local-training-only
lower bound and an unconstrained-communication upper bound (every
publication reaches every compatible receiver each round), both under
the identical evaluation-gated adoption protocol.
